\chapter{Libraries}
\label{chapter:libraries}

This chapter provides documentation for domain-specific or general-purpose libraries developed or integrated into the project. All library modules are defined in CMake using the naming convention \texttt{\textit{Name}\_LIB}. These modules can be linked into your targets via CMake's \texttt{target\_link\_libraries} command. This allows modular inclusion of specific library functionalities based on your build requirements. To use a library, simply add the corresponding library target to your CMake target, for example:
\begin{lstlisting}[style=cppstyle]
target_link_libraries(MyTarget PRIVATE Database_LIB)
\end{lstlisting}




\section{Database}
\label{section:database}

TODO 






\section{Markov}
\label{section:markov}

TODO 









\section{RPG Library}
\label{section:rpg}

The RPG library (\texttt{bin/libs/rpg/}) is an in-development RPG system, written largely for experimentation, inspired by classic World of Warcraft. The library has its own dedicated design document under \texttt{bin/libs/rpg/docs/}, a standalone TeX project that is not included in this manual. Refer to that for the full design.


















\section{Python Plotting Library}\label{section:python_plotting}

The \texttt{Python\_Plotting\_LIB} library provides plotting capabilities for c++ code through the Python plotting module \texttt{PythonPlotter.py}. The library contains the following APIs:
\begin{itemize}
	\item \texttt{PythonPlotter.hpp}, with the Python plotting implementation contained in \texttt{PythonPlotter.py}.
\end{itemize}
The Python plotting module uses \texttt{matplotlib} and \texttt{numpy}, so both must be available to the Python interpreter which loads it. If either package is missing, constructing a \texttt{PythonPlotter} throws an exception when the module fails to import.

\subsection{CMake Setup}

The library is only built when the Python development packages are found (see section \ref{section:python_modules}). The library can be linked to a target using CMake:
\begin{lstlisting}[style=makestyle]
target_link_libraries(MyTarget PRIVATE Python_Plotting_LIB)
\end{lstlisting}

\subsection{Using the PythonPlotter}

The \texttt{PythonPlotter} class provides methods for creating plots. A simple plot can be created by providing an x-axis vector and a y-axis vector:
\begin{lstlisting}[style=cppstyle]
// Sequence to test plotting with.
std::vector<int> x = {0,1,2,3,4,5,6,7};
std::vector<int> y = {0,1,3,6,10,15,21,28};

python_plotting::PythonPlotter plotter;
plotter.plot(x,y);  
\end{lstlisting}
The two vectors must contain the same number of elements. The plotting method is intended to support any numerical vector types. Plot titles and axis labels can be configured using the \texttt{setLabels()} method:
\begin{lstlisting}[style=cppstyle]
plotter.setLabels("My Plot", "X Axis", "Y Axis");
\end{lstlisting}
The configured labels are used by subsequent plots.

The \texttt{PythonPlotter} class can also plot multiple lines using the \texttt{LinesToPlot} type. Each line is represented by a \texttt{LineMetaData} structure containing the data and its plotting properties.
\begin{lstlisting}[style=cppstyle]
// Sequence to test plotting with.
std::vector<int> x = {0,1,2,3,4,5,6,7};

python_plotting::LinesToPlot<int> data = {
	{
		{0,1,3,6,10,15,21,28},
		python_plotting::LineStyle::solid,
		4.0,
		python_plotting::Color::blue
	},
	{
		{0,1,2,4,8,16,32,64},
		python_plotting::LineStyle::dashed,
		1.0,
		python_plotting::Color::red
	}
};

python_plotting::PythonPlotter plotter;
plotter.setLabels("MIATest Plot", "X-Values", "Y-Values");
plotter.plot(x, data);
\end{lstlisting}
Each \texttt{LineMetaData} contains the following properties:
\begin{itemize}\itemsep0em
	\item \texttt{values} -- The data values to plot.
	\item \texttt{style} -- The line style to use.
	\item \texttt{lineWidth} -- The width of the line.
	\item \texttt{color} -- The color of the line.
\end{itemize}
The available line styles are \texttt{solid}, \texttt{dotted}, \texttt{dashed}, and \texttt{dashdot}. Colors can be specified using one of the predefined \texttt{python\_plotting::Color} values:
\begin{lstlisting}[style=cppstyle]
python_plotting::Color::black
python_plotting::Color::red
python_plotting::Color::green
python_plotting::Color::blue
python_plotting::Color::yellow
python_plotting::Color::orange
python_plotting::Color::purple
python_plotting::Color::white
python_plotting::Color::gray
\end{lstlisting}
A hexadecimal color string can also be provided through the \texttt{PlotColor} type:
\begin{lstlisting}[style=cppstyle]
python_plotting::PlotColor color = "#336699";
\end{lstlisting}
The \texttt{PlotColor} type is a variant containing either a predefined \texttt{Color} or a hexadecimal color string.

Verbose output can be enabled using \texttt{setVerboseOutput()}:
\begin{lstlisting}[style=cppstyle]
plotter.setVerboseOutput(true);
\end{lstlisting}
When enabled, the Python plotting module reports information about plotting operations as they occur.

The Python plotting module is installed along with the library according to the configured CMake build mode. This allows the \texttt{PythonPlotter} class to locate and load the Python plotting implementation when it is used.


